
%%======%%======%0%======%%=======%%
\vspace{5mm}
\section{Requisitos do sistema}
\label{sec:Requisitos}
%\addcontentsline{toc}{chapter}{Objetivos}

%%======%%======%0%======%%=======%%

Uma vez bem compreendido o problema a ser sanado e o público alvo, foi necessário listar todos os requisitos do sistema a ser desenvolvido. Para tal, foi feita uma pesquisa em campo, onde o aluno autor se dispôs a conversar com três Professores Pesquisadores, da área de materiais, para melhor compreender as reais necessidades, os professores em questão foram:
\begin{itemize}
	\item Prof. Dr. José R. R. Bortoleto (UNESP);
	\item Prof. Me. Paulo Silas Oliveira (Instituto Federal);
	\item Prof. Dr. Raul Ramos (UNICAMP).
\end{itemize}
As respostas obtidas foram organizadas e sintetizadas a seguir.
\vspace{5mm}
\subsection{Conteúdo esperado para a página inicial}

Quando questionados sobre quais informações deveriam ser apresentadas na página inicial do sistema, foram identificadas três expectativas principais:

\begin{enumerate}
	\item \textbf{Apresentação do laboratório:} disponibilizar uma página introdutória que apresente o laboratório, sua finalidade, suas áreas de atuação e suas principais características;
	
	\item \textbf{Pesquisas e frentes de trabalho:} apresentar as principais linhas de pesquisa, projetos e frentes de trabalho desenvolvidas pelo grupo;
	
	\item \textbf{Panorama das atividades do grupo:} disponibilizar uma síntese das atividades realizadas, contemplando acontecimentos recentes, projetos em andamento, resultados e demais novidades relevantes.
\end{enumerate}

Dessa forma, observa-se que a página inicial é percebida como um meio de divulgação e contextualização das atividades desenvolvidas pelo laboratório.
\vspace{5mm}
\subsection{Funcionalidades esperadas para o sistema}

Em relação às funcionalidades esperadas, as respostas indicaram a necessidade de uma plataforma que além de apresentar informações institucionais também ofereça recursos para integração, organização e gerenciamento das atividades do grupo. Entre as principais funcionalidades mencionadas, destacam-se:

\begin{enumerate}
	\item \textbf{Integração entre membros e público externo:} possibilitar a comunicação e a aproximação entre os integrantes do grupo e pessoas externas ao laboratório, facilitando o acesso às informações e aos trabalhos desenvolvidos;
	
	\item \textbf{Agendamento e gerenciamento de recursos:} permitir a realização de agendamentos, a consulta de disponibilidade e a organização do uso dos recursos e equipamentos do laboratório;
	
	\item \textbf{Repositório de trabalhos:} disponibilizar um espaço destinado ao armazenamento e à consulta de trabalhos acadêmicos, projetos, publicações e demais produções desenvolvidas pelo grupo;
	
	\item \textbf{Gerenciamento de reagentes e insumos:} possibilitar a consulta às informações relacionadas aos reagentes e materiais disponíveis no laboratório;
	
	\item \textbf{Informações sobre máquinas e equipamentos:} apresentar os equipamentos disponíveis, suas características e, quando aplicável, sua disponibilidade para utilização;
	
	\item \textbf{Divulgação de novidades:} disponibilizar informações sobre acontecimentos recentes, atividades, projetos, resultados e demais novidades relacionadas ao grupo.
\end{enumerate}

A partir dessas respostas, verifica-se uma expectativa de que o sistema atue como uma plataforma integrada de informações e gerenciamento, reunindo em um único ambiente tanto os aspectos institucionais e acadêmicos do laboratório quanto os recursos necessários à organização de suas atividades, tais dados foram abstraídos pelo aluno autor, e pode ser desenvolvida a seguinte lista de requisitos ligados a parte técnica do projeto:
\begin{itemize}
	\item ser responsivo;
	\item possuir contraste adequado;
	\item utilizar estrutura HTML semântica;
	\item funcionar nos navegadores exigidos;
	\item apresentar desempenho adequado
	\item ser simples de usar.
	\item apresentar informações do laboratório;
	\item disponibilizar trabalhos;
	\item permitir agendamento;
\end{itemize}
